%%%%%%%%%%%%%%%%%%%%%%%%%%%%%%%%%%%%%%%%%
% SINGLE-LINE COPY FOR BULLET LISTS
%
% Why: PDF text extraction is line-based. When a bullet wraps onto two
% visual lines, every extractor hands the text back with a hard line break
% in the middle, so pasting a job description into a text editor (or into
% an ATS field) gives one paragraph per VISUAL line instead of one per
% bullet. That is the difference between a readable paragraph and a wall
% of half-sentences.
%
% How: PDF 32000-1 14.9.3 lets a content stream mark a span of text and
% declare what it actually says (/ActualText). Extractors that honour the
% span replace the whole marked run by that string, so a wrapped bullet is
% copied back as one continuous line. This file puts every item of every
% itemize list inside such a span, with the item's own source text as the
% /ActualText.
%
% The visible PDF does not change: the span only adds marked content to the
% content stream. The bullet label, the fonts, the line breaking and the
% vertical spacing are exactly what LaTeX would have produced -- the
% ORIGINAL \item is still the one that runs, so the list layout, the
% paragraph setup and the label placement are untouched.
%
% Note on the readers: Poppler (pdftotext), PDFium (Chrome, Brave) and
% Acrobat honour /ActualText. macOS Preview does not -- PDFKit copies
% strictly what is drawn, line by line, and ignores marked content. There
% is no way to change that from inside the PDF without turning the text
% into outlines, which would destroy the selectable-text requirement of
% the ATS layout.
%%%%%%%%%%%%%%%%%%%%%%%%%%%%%%%%%%%%%%%%%

\RequirePackage{accsupp}

\makeatletter

%----------------------------------------------------------------------------------------
%	BRANCH SELECTION
%
%	\at@ifempty{<text>}{<empty code>}{<non-empty code>}
%
%	The branch is SELECTED inside the \if but EXPANDED after the \fi, on
%	purpose: if code that runs later were still inside the conditional,
%	the pending \fi would be captured as part of the text that follows.
%----------------------------------------------------------------------------------------
\def\at@empty{}
\def\at@branchA#1#2{#1}%
\def\at@branchB#1#2{#2}%
\long\def\at@ifempty#1#2#3{%
	\edef\at@tmp{\detokenize{#1}}%
	\ifx\at@tmp\at@empty
		\let\at@branch\at@branchA
	\else
		\let\at@branch\at@branchB
	\fi
	\at@branch{#2}{#3}%
}

%----------------------------------------------------------------------------------------
%	EMIT ONE ITEM
%
%	The item is started the way LaTeX itself starts an item that has no
%	custom label: \@noitemargtrue\@item[\@itemlabel]. Writing it that way
%	instead of \item avoids \@ifnextchar, whose lookahead can mistake the
%	text of an item for the optional argument of the next one. \reserved@a
%	is cleared around the call because that is the scratch register the
%	lookahead writes into.
%
%	\leavevmode starts the paragraph BEFORE the marked-content span is
%	opened, so the opening operator (BDC) becomes part of the paragraph
%	instead of a page-level literal. A marked-content span has to be
%	balanced inside a single page's content stream, so a bullet that broke
%	across two pages would leave its opening operator on one page and its
%	closing one on the next -- malformed, and poppler says so ("Mismatched
%	EMC operator"). The penalties below keep the bullet together on one
%	page, which is also how a bullet that was already whole on the page
%	would have been set.
%
%	The text of the item is captured as it is written: the space that
%	follows \item is consumed by TeX's tokenizer, so the copied text
%	starts with the first word of the bullet.
%----------------------------------------------------------------------------------------
\long\def\at@emititem#1{%
	\let\reserved@a\@empty
	\@noitemargtrue\@item[\@itemlabel]%
	\let\reserved@a\@empty
	\interlinepenalty=10000
	\clubpenalty=10000
	\widowpenalty=10000
	\leavevmode
	\BeginAccSupp{method=pdfstringdef,unicode=true,ActualText={#1}}%
		#1%
	\EndAccSupp{}%
}

%----------------------------------------------------------------------------------------
%	SPLIT THE CAPTURED BODY ON \item
%
%	The body of a list is captured whole (see \at@grabbody) and then walked
%	item by item here. The body ends with "\item\item": the second \item
%	closes the last scan, which yields an empty chunk, and the empty chunk
%	is what stops the walk.
%----------------------------------------------------------------------------------------
\def\at@split#1\item{%
	\at@ifempty{#1}{}{\at@emititem{#1}\at@split}%
}

%----------------------------------------------------------------------------------------
%	CAPTURE THE BODY
%
%	\at@grabbody grabs everything between \begin{itemize} and \end, drops the
%	leading \item (the marker of the first item, which is re-created by
%	\at@emititem) and appends the terminator that makes the scans stop.
%	The \end is put back so the environment closes as usual.
%----------------------------------------------------------------------------------------
\long\def\at@stripitem#1\item#2\@nil{%
	\expandafter\def\expandafter\at@rest\expandafter{#2\item\item}%
}
\long\def\at@grabbody#1\end{%
	\gdef\at@body{#1}%
	\at@ifempty{#1}{}{%
		\expandafter\at@stripitem\at@body\@nil
		\expandafter\at@split\at@rest
	}%
	\end%
}

%----------------------------------------------------------------------------------------
%	HOOK
%
%	\itemize is wrapped, not replaced: the original starts the list, then
%	the body is captured and re-emitted item by item. A list without items
%	is left alone (nothing to mark).
%----------------------------------------------------------------------------------------
\let\at@origitemize\itemize
\renewcommand{\itemize}{\at@origitemize\at@grabbody}

\makeatother
